\documentclass[12pt]{article}
\usepackage{amsmath,amssymb,amsthm}
\usepackage{geometry}
\usepackage{hyperref}
\usepackage{booktabs}
\usepackage{enumitem}
\usepackage{parskip}
\geometry{letterpaper, margin=1in}

\newtheorem{theorem}{Theorem}[section]
\newtheorem{definition}[theorem]{Definition}
\newtheorem{proposition}[theorem]{Proposition}
\newtheorem{conjecture}[theorem]{Conjecture}
\newtheorem{corollary}[theorem]{Corollary}
\newtheorem{remark}[theorem]{Remark}

\newcommand{\lP}{l_P}
\newcommand{\EP}{E_P}
\newcommand{\tP}{t_P}
\newcommand{\mP}{m_P}
\newcommand{\SSV}{\text{SSV}}
\newcommand{\PSR}{\text{PSR}}
\newcommand{\nuP}{\nu_P}
\newcommand{\nuC}{\nu_C}
\newcommand{\alphageom}{\alpha_{\rm geom}}
\newcommand{\lambdabar}{\bar{\lambda}}

\title{\textbf{Inertial Mass from Zitterbewegung:\\
A Two-Level Resonance Hierarchy in Conscious Point Physics}\\[6pt]
{\large Companion Paper to ``Mechanistic Derivation of Relativistic Effects\\
via Space Stress Vector (SSV) in the Dipole Sea'' (Version~16)}}

\author{Thomas Lee Abshier, ND\\
Hyperphysics Institute\\
\texttt{https://hyperphysics.com}\\
\texttt{drthomas007@protonmail.com}}

\date{13 March 2026 --- Version~2}

\begin{document}
\maketitle

\noindent\textbf{Keywords:} Conscious Point Physics, inertial mass,
Zitterbewegung, Compton frequency, Planck action, Dipole Sea polarization,
two-level resonance hierarchy, $E = mc^2$, eDP standing wave, non-radiation.

\begin{abstract}
We show that inertial mass in Conscious Point Physics (CPP) arises from a
two-level Zitterbewegung (ZBW) resonance hierarchy in the polarized Dipole
Sea.  At the Planck scale (Level~1), each elementary Dipole Particle (DP)
oscillates at frequency $\nuP = 1/(2\tP)$, carrying action $\hbar =
\EP\tP$ per half-cycle.  At the particle scale (Level~2), the coherent eDP
polarization cloud surrounding an unpaired electron-type Conscious Point
oscillates collectively as a radial standing wave at the Compton frequency
$\nuC = mc^2/\hbar$, reflecting between the nucleation center and the
thermal boundary of the cloud.  The particle's rest mass is the energy of
one complete Compton resonance cycle, $mc^2 = \hbar\nuC$, derived here from
standing-wave boundary conditions rather than assumed.  The ratio of scales
$\nuP/\nuC = \mP/(2m)$ shows that one Compton cycle spans $\sim\mP/m$
Planck ticks, realizing the assembly principle that larger organizational
units are built from smaller ones.  The relation $E = mc^2$ follows from
the Absolute Moment step-budget constraint.  The classical non-radiation of
ground-state electrons is explained as a driven steady state maintained by
the eternal nucleation signal of the unpaired eCP.  All four DP types
participate equally in the polarization cloud.  No new postulates are
introduced.
\end{abstract}

\section{Introduction}
\label{sec:intro}

The preceding companions established:
(i) the universal Planck tick $\tP$ as necessary for DI-bit conservation
(Absolute Moment companion);
(ii) the Dipole Sea stiffness $C = \alphageom\EP/\lP^3$ from 600-cell
Voronoi geometry (Stiffness~$C$ companion);
(iii) $\hbar = \EP\tP$ as the action per Planck ZBW half-cycle, with
mechanistic accounts of superposition, measurement, and collapse (Born-rule
companion).

The present companion closes the quantum-to-classical bridge: inertial mass
is not a primitive property of Conscious Points but the energy stored in a
two-level ZBW resonance of the surrounding eDP polarization cloud.  The
central result is
\begin{equation}
  mc^2 \;=\; \hbar\nuC, \qquad \nuC = \frac{mc^2}{\hbar},
  \label{eq:mass-compton}
\end{equation}
where $\nuC$ is the Compton frequency.  The left side is the measured rest
energy; the right side identifies $\nuC$ as the frequency of a radial
standing wave derived from CPP boundary conditions.

\begin{remark}[Relation to Schr\"odinger's ZBW]
Schr\"odinger (1930) identified a rapid oscillatory motion of the Dirac
electron at frequency $2mc^2/\hbar$ (twice the Compton frequency), arising
from interference between positive- and negative-energy spinor components.
The CPP picture gives this a concrete mechanism: the factor of 2 reflects
the two-way radial transit of the standing wave (outward leg and return
leg).  CPP therefore explains what Dirac's equation describes.
\end{remark}

\section{The Two-Level ZBW Hierarchy}
\label{sec:hierarchy}

\begin{definition}[Level 1 --- Planck ZBW]
Every elementary DP oscillates between attraction to its pair-mate and
attraction to the surrounding DP sea at the Absolute Moment frequency:
\begin{equation}
  \nuP \;=\; \frac{1}{2\tP}.
  \label{eq:nu-planck}
\end{equation}
The action per half-cycle is $\hbar = \EP\tP$.  This is the elementary
unit of energetic organization in the CPP lattice.
\end{definition}

\begin{definition}[Level 2 --- Compton ZBW]
The coherent eDP polarization cloud surrounding an unpaired eCP oscillates
collectively as a radial standing wave at the Compton frequency
\begin{equation}
  \nuC \;=\; \frac{mc^2}{\hbar},
  \label{eq:nu-compton}
\end{equation}
propagating outward from the eCP nucleation center, reflecting at the
thermal boundary, returning to center, and repeating.
\end{definition}

The ratio of the two frequencies is a pure mass ratio:
\begin{equation}
  \frac{\nuP}{\nuC} \;=\; \frac{\mP}{2m},
  \label{eq:ratio}
\end{equation}
where $\mP = \sqrt{\hbar c/G}$ is the Planck mass.  One Compton cycle spans
$N_{\rm Planck} = \mP/(2m)$ Planck ticks.

\begin{center}
\renewcommand{\arraystretch}{1.4}
\begin{tabular}{@{}llll@{}}
\toprule
Level & Entity & Frequency & Energy per cycle \\
\midrule
1 (Planck) & Individual DP & $1/(2\tP)$ & $\EP = \hbar/\tP$ \\
2 (Compton) & eDP cloud & $mc^2/\hbar$ & $mc^2 = \hbar\nuC$ \\
\bottomrule
\end{tabular}
\end{center}

For the electron: $N_{\rm Planck} = \mP/(2m_e) \approx 1.2\times10^{22}$.
One electron Compton cycle is a coherent aggregate of $\sim10^{22}$ Planck
ZBW ticks --- the CPP realization of the assembly principle.

\section{The eDP Polarization Cloud}
\label{sec:cloud}

\subsection{Equal participation of all four DP types}

The central unpaired eCP emits the same elementary charge signal $\pm1$ to
all surrounding DPs regardless of their type (Version~16, \S2; Stiffness~$C$
companion, \S2.2).  No differential force distinguishes eDP, qDP, or hybrid
DP types at the CP charge level.  All four types therefore participate
equally in the polarization cloud.

This is confirmed by the near-$c$ limiting argument: at relativistic
velocities the PSR contracts severely, leaving no time budget for type
segregation.  Since no segregation occurs even in this extreme case, it
cannot occur at any velocity.  The energy accounting is uniform: each DP
in the cloud contributes one Planck ZBW unit regardless of type.

\subsection{Standing-wave boundary conditions}

Two boundaries define the eDP cloud:

\begin{enumerate}[noitemsep]
  \item \emph{Inner boundary (nucleation center):} The unpaired eCP provides
    a persistent SSV source that reflects the inward-propagating SSV
    wavefront via the CP Exclusion Rule (no second CP can occupy the same
    Grid Point at the same $t_{\rm abs}$).

  \item \emph{Outer boundary (thermal radius $r_{\rm th}$):} Beyond $r_{\rm
    th}$, incoherent thermal SSV from Brownian collisions and solitonic
    superpositions dominates the nucleation signal, randomizing DP
    orientations faster than they can be re-aligned.
\end{enumerate}

The lowest-mode standing wave fits one half-wavelength between center and
thermal boundary:
\begin{equation}
  r_{\rm th} \;=\; \frac{c}{2\nuC} \;=\; \frac{\hbar}{2mc}
  \;=\; \tfrac{1}{2}\lambdabar_C,
  \label{eq:rth}
\end{equation}
where $\lambdabar_C = \hbar/(mc)$ is the reduced Compton wavelength.
Equation~\eqref{eq:rth} identifies the thermal radius of the polarization
cloud with the reduced Compton wavelength --- a result known empirically
in quantum mechanics, here derived from standing-wave boundary conditions.

\subsection{Why ground-state particles do not radiate}

The nucleation center (unpaired eCP) is eternal: it emits its SSV signal
at every Absolute Moment tick, without diminution, for as long as it exists.
The eDP cloud is therefore a \emph{driven steady state}, continuously
re-nucleated against Brownian disruption.  There is no lower-energy
configuration available, because the nucleation signal is permanent and the
CP Exclusion Rule prevents the cloud from collapsing to zero radius.

This is the CPP resolution of the classical radiation paradox: ground-state
electrons do not radiate because their polarization clouds are driven steady
states, not metastable configurations seeking a lower energy level.  The
``ground state'' is the resonant radius~\eqref{eq:rth} at which the
nucleation rate exactly balances the Brownian disruption rate.

\section{Rest Mass as Standing-Wave Energy}
\label{sec:mass}

\begin{proposition}[Rest mass from Compton standing wave]
\label{prop:mass}
The rest mass of a massive particle equals the energy of one complete
Compton radial oscillation of its eDP polarization cloud:
\begin{equation}
  mc^2 \;=\; \hbar\nuC.
  \label{eq:mc2}
\end{equation}
\end{proposition}

\begin{proof}
The standing wave completes one round trip in period $T_C = 2r_{\rm
th}/c$.  The energy stored in the cloud is the product of the Planck ZBW
action quantum $\hbar$ and the number of Planck ticks per Compton cycle:
\begin{equation}
  E_{\rm cloud} \;=\; N_{\rm Planck}\cdot\frac{\hbar}{2\tP}
  \;=\; \frac{\nuP}{\nuC}\cdot\frac{\hbar}{2\tP}
  \;=\; \frac{1}{2\tP\nuC}\cdot\frac{\hbar}{2\tP}\cdot\frac{1}{2\tP}
  \;=\; \hbar\nuC,
  \label{eq:E-cloud}
\end{equation}
where we used $\nuP = 1/(2\tP)$.  Setting $E_{\rm cloud} = mc^2$ gives
$mc^2 = \hbar\nuC$ and determines $\nuC$ consistently.
\end{proof}

\begin{remark}[Unified energy accounting]
This result unifies energy across all scales within a single ZBW unit-count
framework:
\begin{itemize}[noitemsep]
  \item \emph{Rest mass:} $\mP/(2m)$ Planck ZBW units organized into one
    Compton standing wave.
  \item \emph{Kinetic energy:} Additional Planck ZBW units organized by
    acceleration (dynamic polarization of the moving cloud).
  \item \emph{Photon energy $h\nu$:} ZBW units organized into a propagating
    eDP wavefront at frequency $\nu$ (propagating eDP wavefront with no
    nucleation center), hence no rest mass and no standing-wave boundary
    condition.
\end{itemize}
The prior CPP definition of energy as ``order of the Dipole Sea, measured
by DP orientation and separation'' is the macroscopic description of this
microscopic ZBW unit accounting.
\end{remark}

\section{\boldmath$E = mc^2$ from the Absolute Moment Step Budget}
\label{sec:emc2}

The step-budget relation (Absolute Moment companion, \S5.1) is:
\begin{equation}
  \lP^2 \;=\; (c\,\Delta\tau)^2 + |\mathbf{d}_{\rm spatial}|^2.
  \label{eq:budget}
\end{equation}
For a particle at rest, $|\mathbf{d}_{\rm spatial}|=0$ and the full budget
goes to the timelike direction.  For a particle moving at speed $v$:
\begin{equation}
  \Delta\tau \;=\; \tP\sqrt{1 - v^2/c^2},
  \qquad
  |\mathbf{d}_{\rm spatial}| \;=\; \lP\,\frac{v}{c}.
  \label{eq:budget-moving}
\end{equation}
The total relativistic energy follows from Proposition~\ref{prop:mass} and
the Lorentz factor implied by~\eqref{eq:budget-moving}:
\begin{equation}
  E_{\rm total} \;=\; \frac{mc^2}{\sqrt{1-v^2/c^2}} \;=\; \gamma\,mc^2.
  \label{eq:relativistic-energy}
\end{equation}
The conversion factor $c^2$ in $E = mc^2$ is geometric: it is
$(\lP/\tP)^2$, the square of the lattice step-to-tick ratio --- a property
of the 600-cell, not an independent empirical input.

\section{Consistency with Previous Companions}
\label{sec:consistency}

Every element of the present derivation uses only machinery already
established:

\begin{itemize}[noitemsep]
  \item The \emph{Absolute Moment tick} $\tP$ fixes $\hbar = \EP\tP$ and
    $\nuP = 1/(2\tP)$ (Absolute Moment companion).
  \item The \emph{stiffness $C$} governs the elastic response of the eDP
    cloud to SSV perturbation and enters the PSR formula of Version~16.
  \item The \emph{CP Exclusion Rule} provides the inner reflection boundary
    that creates the standing wave.
  \item The \emph{shell-broadcast mechanism} at speed $c$ drives the
    outward-propagating leg of the standing wave.
  \item The \emph{step-budget constraint}~\eqref{eq:budget} reproduces
    $E=mc^2$ and relativistic kinematics (Absolute Moment companion, \S5.1).
\end{itemize}

No element contradicts the SR derivation of Version~16 or the three
preceding companions.

\section{Open Problems}
\label{sec:open}

\textbf{(1) Numerical verification of $N_{\rm Planck}$ for known particles.}
The ratio $\mP/(2m)$ gives $N_{\rm Planck}\approx1.2\times10^{22}$ for
the electron, $\approx2.2\times10^{19}$ for the proton.  A Monte-Carlo
extension of the existing geometric-strain routine (Version~16, A.8.1)
could verify that a coherent aggregate of this many Planck ZBW ticks
reproduces the correct Compton wavelength.

\textbf{(2) Coupling to general relativity.}  The eDP cloud's energy
density sources the SSV gradient that curves the effective 600-cell
geometry, producing gravitational effects.  This is the subject of the next
companion.

\textbf{(3) Higher resonance modes.}  The fundamental (half-wavelength)
mode is identified with rest mass.  Higher radial modes ($n =
1,2,\ldots$) may correspond to excited states or internal structure of
composite particles.  Their quantization condition is deferred.

\textbf{(4) Electron ground state in hydrogen.}  The minimum orbital radius
below which no stable 12-edge selection sequence closes is a candidate for
the Bohr radius, derivable from the 600-cell geometry (noted in the
Born-rule companion, \S7).

\section{Conclusion}
\label{sec:conclusion}

Inertial mass in CPP is the energy stored in a two-level ZBW resonance
hierarchy:

\begin{enumerate}
  \item \textbf{Level 1 (Planck ZBW):} Each elementary DP carries action
    $\hbar = \EP\tP$ per half-cycle at frequency $1/(2\tP)$.

  \item \textbf{Level 2 (Compton ZBW):} The coherent eDP cloud surrounding
    an unpaired eCP oscillates as a radial standing wave at $\nuC =
    mc^2/\hbar$, with thermal radius $r_{\rm th} = \hbar/(2mc)$
    (the reduced Compton wavelength).

  \item \textbf{Rest mass $= \hbar\nuC$:} Derived from the standing-wave
    energy of the Compton oscillation, not assumed.

  \item \textbf{Non-radiation:} Ground-state particles do not radiate
    because the eDP cloud is a driven steady state maintained by the
    eternal nucleation signal of the unpaired eCP.

  \item \textbf{$E = mc^2$:} Follows from the Absolute Moment step-budget
    constraint; the factor $c^2 = (\lP/\tP)^2$ is a geometric property
    of the 600-cell lattice.
\end{enumerate}

This completes the quantum-to-classical bridge.  The same 600-cell lattice,
PCD cycle, and SSV broadcast that drive special relativity and
electromagnetism also drive inertial mass, with no additional postulates.
The coupling of this mass to gravitational curvature is the subject of the
next companion.

\section*{References}

\begin{enumerate}[label={[\arabic*]}]
  \item Abshier, T.~L. (2026). Mechanistic Derivation of Relativistic
    Effects via Space Stress Vector (SSV) in the Dipole Sea.  Version~16.

  \item Abshier, T.~L. (2026). The Absolute Moment Postulate: Necessity,
    Consistency, and the PCD Cycle in CPP.  Version~3 (companion 1).

  \item Abshier, T.~L. (2026). Microscopic Origin of the Dipole Sea
    Stiffness $C$ in CPP.  Version~2 (companion 2).

  \item Abshier, T.~L. (2026). Quantum Probability and the Classical
    Transition in CPP: A Mechanistic Framework.  Version~2 (companion 3).

  \item Abshier, T.~L. (2025--2026). Conscious Point Physics Series.
    \url{https://hyperphysics.com/papers}

  \item Schr\"odinger, E. (1930). \"Uber die kr\"aftefreie Bewegung in der
    relativistischen Quantenmechanik. \textit{Sitzungsber.\ Preuss.\ Akad.\
    Wiss.}, Phys.-Math.\ Kl., 418--428.
\end{enumerate}

\bigskip
\noindent\textbf{GitHub (geometric-strain routine, basis for future
ZBW hierarchy Monte-Carlo extension):}\\
\url{https://github.com/tlabshier/CPP/blob/main/600-cell_special-relativity_emergence/600cell_monte_carlo_voronoi_k_fit.py}

\end{document}
